%! Planning a Study — Unit 3, Lesson 3.1 — Guided Notes & Practice
% THE in-class document, in the MAIN IDEAS / NOTES shape (see LESSON_SHAPE.md §1), at the
% COMPACT budget: the vocabulary box, then ONE guidednotes table of FOUR rows (three instruction
% rows + Guided Practice), 14 problems, 2 pages.
%   I DO   : the teacher fills the definition lines and works the FIRST problem of each grid
%            (1, 4, 7).
%   WE DO  : the class works the rest; the last row, Guided Practice, is worked together.
%   YOU DO : nothing on this page — ../homework/ IS the individual practice.
% THE TRAP is problem 2 (the coach's varsity players sound reasonable, but the coach chose them).
% THE CRUX is problem 10, in the last instruction row: "the sleep lab measured everything so
% carefully that it is an experiment" — careful measurement is not a treatment (DAT-2.A.3/4).
%
% THE CONTEXTS (instruction, guided practice, and homework each use a different one)
%  Row 1 — three one-line methods: radio call-in, a coach's varsity players, 60 random IDs.
%  Row 2 — a principal's homework-hours survey: 100 of the school's 1,200 students at random.
%  Row 3 — four one-line studies: coffee records, paper vs laptop quizzes, music and GPA, a
%          sleep lab's wrist monitors (the crux).
%  Row 4 (GP) — a university's 20,000 students: a random survey of 250, then a flashcard-app
%          study with random assignment.
%
% PAGE LOCKSTEP with notes_key/: \termblank -> \termans, \blank -> \ans, 4th arg of \pcell.
% TABLE MECHANICS: inside a cell `\\` ENDS THE ROW — use \par. Long rows are split into
% sub-rows (a bare `\\` then `& ...`) so the table can break between them.
% NO SKETCH QUESTIONS: nothing is drawn by the student.
\documentclass[12pt]{article}
\usepackage{apstats-article}
\usepackage{apstats-boxes}

\begin{document}

\pageheader{Unit 3, Lesson 3.1}{Guided Notes \& Practice}
% NAMESTRIP: no \namedateperiod here — the name row lives on the cover only.

\vspace{2pt}

\boxguard[16]
\begin{vocabbox}
\small
Fill in each term \textbf{as we name it} in the notes below.
\par
\termblank{Population}
\par
\termblank{Sample}
\par
\termblank{Observational study}
\par
\termblank{Experiment}
\par
\end{vocabbox}

\small
\begin{guidednotes}
% ── ROW 1: who decided who is in the data? ───────────────────────────────────
\mainidea[Before you trust it]{Ask how} &
Before you trust a number, ask \emph{how the data were collected}. In a
\textbf{chance-based method}, a \blank{1.6cm} process decides who is in the data. In a
\textbf{non-chance method}, \blank{1.4cm} decide --- the researcher takes whoever is convenient,
or the respondents volunteer. Methods that do not rely on chance give \blank{2.8cm}
conclusions.
\notesprompt{Who decided who is in the data: chance or people?}
\begin{probgrid}{|Y|Y|Y|}
\pcell{1}{A radio host asks listeners to call in and vote on a new stadium.}{1.1cm}{} &
\pcell{2}{A coach asks the varsity players if the school should spend more on sports.}{1.1cm}{} &
\pcell{3}{A computer randomly picks 60 ID numbers from the full student roster.}{1.1cm}{} \\ \hline
\end{probgrid}
\\ \hline
% ── ROW 2: who we want to know about vs who we measure ───────────────────────
\mainidea[Who we ask]{Sample} &
The \textbf{population} is \blank{0.8cm} the individuals we want to know about; a
\textbf{sample} is the \blank{1.4cm} of it we actually measure; a \textbf{census} measures
\blank{1.2cm} individual. Generalize only to the \blank{2.2cm} the sample came from, and only
when chance chose it.\par\vspace{3pt}
A principal asks about weekly homework hours. A computer picks \textbf{100} of her school's
\textbf{1{,}200} students at random from the roster.
\notesprompt{Name the parts of her study.}
\begin{probgrid}{|Y|Y|Y|}
\pcell{4}{Name the population and the sample.}{1.4cm}{} &
\pcell{5}{Why might she not take a census?}{1.4cm}{} &
\pcell{6}{Does her result describe all high schoolers in the state?}{1.4cm}{} \\ \hline
\end{probgrid}
\\ \hline
% ── ROW 3: imposing a treatment — THE CRUX ───────────────────────────────────
\mainidea[Watch or do]{Study type} &
In an \textbf{observational study}, researchers observe and measure but do not
\blank{1.4cm} a treatment; a sample survey is one kind. In an \textbf{experiment},
researchers \blank{1.4cm} treatments to the individuals and measure the response. Only a
well-designed \blank{2.0cm} can show cause and effect; an observational study shows only an
\blank{2.2cm}.
\\
& \notesprompt{Observational study or experiment? Say why.}
\begin{probgrid}{|Y|Y|}
\pcell{7}{Researchers read the medical records of 2{,}000 adults and compare heart disease in coffee drinkers and non-drinkers.}{1.5cm}{} &
\pcell{8}{A teacher randomly assigns half her classes to take quizzes on paper and half on laptops, then compares scores.}{1.5cm}{} \\ \hline
\pcell{9}{Of 300 teens surveyed, those who play an instrument have higher GPAs. Does playing an instrument raise GPA?}{1.5cm}{} &
\pcell{10}{A sleep lab fits 50 adults with wrist monitors and records their exact sleep and caffeine for a month. ``It was measured so carefully that it is an experiment.'' Agree?}{1.5cm}{} \\ \hline
\end{probgrid}
\\ \hline
% ── ROW 4: guided practice — WE DO, together, fresh context ──────────────────
\mainidea[Guided practice]{The Study App} &
A university wants to know about study habits of \emph{all} \textbf{20{,}000} of its students.
\textbf{Study A:} a computer picks 250 students at random from the registrar's list, and an app logs each one's
weekly study hours to the minute. \textbf{Study B:} 120 intro-biology students are randomly assigned
to use a flashcard app or not, and their exam scores are compared.
\notesprompt{We work these together.}
\begin{probgrid}{|Y|Y|}
\pcell{11}{Study A: name the population and the sample.}{1.5cm}{} &
\pcell{12}{Is Study A chance-based? Can its result describe all 20{,}000 students?}{1.5cm}{} \\ \hline
\pcell{13}{Study A's hours were logged to the minute. Does that make it an experiment?}{1.5cm}{} &
\pcell{14}{Is Study B observational or an experiment? What could it show that A cannot?}{1.5cm}{} \\ \hline
\end{probgrid}
\\ \hline
\end{guidednotes}

% ── THE NOTES END AT GUIDED PRACTICE ─────────────────────────────────────────
% There is deliberately no "you do" here: ../homework/ is the individual practice,
% started in class in the last 8 minutes while the teacher circulates.

\end{document}
